\documentclass[11pt]{article}
\usepackage[a4paper,margin=27mm]{geometry}
\usepackage[T1]{fontenc}
\usepackage[utf8]{inputenc}
\usepackage{lmodern,microtype}
\usepackage{amsmath,amssymb,amsthm,mathtools,bm}
\usepackage{booktabs,array,enumitem,float}
\usepackage{xcolor}
\usepackage[colorlinks=true,linkcolor=blue!45!black,citecolor=blue!45!black,urlcolor=blue!55!black]{hyperref}
\usepackage[nameinlink,noabbrev]{cleveref}
\usepackage[backend=biber,style=numeric,sorting=nyt,maxbibnames=99]{biblatex}
\addbibresource{references.bib}
\AtBeginBibliography{\sloppy}
\setlength{\emergencystretch}{2em}

\newtheorem{theorem}{Theorem}[section]
\newtheorem{lemma}[theorem]{Lemma}
\newtheorem{proposition}[theorem]{Proposition}
\newtheorem{corollary}[theorem]{Corollary}
\newtheorem{remark}[theorem]{Remark}
\newcommand{\E}{\mathbb E}
\newcommand{\R}{\mathbb R}
\newcommand{\C}{\mathbb C}
\newcommand{\Hh}{\mathbb H}
\newcommand{\tr}{\operatorname{tr}}
\newcommand{\Herm}{\operatorname{Herm}}
\newcommand{\Var}{\operatorname{Var}}
\newcommand{\Gr}{\operatorname{Gr}}
\newcommand{\cO}{\mathcal O}
\newcommand{\cL}{\mathcal L}
\newcommand{\ip}[2]{\langle #1,#2\rangle}
\newcommand{\norm}[1]{\lVert #1\rVert}

\hypersetup{
 pdftitle={All-Field Morse--Bott Stability at Critical Homogeneous Orbits: Half-Grassmann No-Spinodal Rigidity and Exact Jacobi Metastability},
 pdfauthor={Lluis Eriksson},
 pdfkeywords={positive-definite kernel, homogeneous orbit, Morse--Bott, Grassmannian, matrix Bingham, Jacobi ensemble, metastability}
}

\title{\bfseries All-Field Morse--Bott Stability at Critical Homogeneous Orbits\\[3pt]
\large Half-Grassmann No-Spinodal Rigidity and Exact Jacobi Metastability}
\author{Lluis Eriksson}
\date{August 15, 2026}

\begin{document}
\maketitle

\begin{abstract}
Let a compact group act orthogonally on a sphere and let $\mu$ be the
invariant law of a proper antipodal orbit.  We study the orbital potential
$U_t(a)=\int\exp(t\langle a,u\rangle)\,d\mu(u)$.  At a point of the source
orbit, every spherical-harmonic coefficient of the orbit average is a
squared norm.  Combining this positivity with the strictly positive
Gegenbauer--Bessel coefficients of the exponential kernel gives an exact
negative spherical-Laplacian identity at every field $t>0$.  Provided the
source orbit is critical---automatically so when the normal isotropy
representation has no fixed vector---irreducibility, or a transitive symmetry
of its irreducible normal blocks, upgrades this trace identity to strict
Morse--Bott maximality for every field.  We give an explicit torus-orbit
counterexample showing why the criticality hypothesis cannot be omitted.

Applied to the centered projector embeddings of the real, complex, and
quaternionic half-Grassmannians, the theorem proves all-field local rigidity
of the balanced matrix--Bingham branch and excludes any radial spinodal.  In
the complex case we go further.  For every two-block external spectrum we
derive exact constrained-Hessian eigenvalues from the reversible Jacobi
generator.  Their low-field factors show a sharp metastability threshold:
unbalanced strata with positive multiplicity $k>2r/3$ are genuine local
maxima at weak field, whereas those with $k\le 2r/3$ have an unstable block
split.  A Stein inequality proves the required Hessian sign in all
sufficiently high Taylor degrees and leaves an explicit finite low-degree
gate.  Thus a global competing phase, if it exists, must arise by nonlocal
coexistence rather than loss of stability of the balanced branch.  We do not
claim global spectral optimality or the resulting all-distortion
rate--distortion function.
\end{abstract}

\section{Introduction}

Matrix--Bingham normalizers on Grassmannians are orbital Laplace transforms
and hypergeometric functions of matrix argument.  Their evaluation,
asymptotics, and information geometry are well developed
\cite{James1964,Muirhead1982,Kent1994,BagyanRichards2024,CazzellaEtAl2026}.
A different problem asks which external spectrum maximizes such a normalizer
under trace and Frobenius constraints.  The constraints do not define a
majorization chain, so standard Schur or orbital-integral monotonicity does
not decide the question \cite{ItzyksonZuber1980,Sra2016,McSwiggenSahi2026}.

The rank-two case $\Gr_{\C}(2,4)$ admits an exact all-field solution
\cite{ErikssonRankTwo2026}; arbitrary rank has previously been controlled at
high field, yielding an exact high-fidelity rate--distortion segment
\cite{ErikssonHighFidelity2026}.  The missing finite-field issue is global:
can an intermediate external spectrum overtake the balanced two-block field?
This paper does not resolve that global ordering.  It instead closes its
local part at every field and proves that a seemingly natural saddle-based
shortcut is false.

The first observation is representation-free.  Classical Schoenberg and
Gangolli theory characterizes positive-definite zonal kernels by nonnegative
harmonic coefficients \cite{Schoenberg1942,Gangolli1967}.  If the center of
the orbital potential lies on the same orbit as the integration measure, the
orbital harmonic coefficient is itself a squared mean harmonic.  The
spherical Laplacian is consequently strictly negative.  Under an explicit
criticality condition and a normal-isotropy hypothesis this trace sign becomes
a complete normal-Hessian sign, giving an all-field Morse--Bott theorem.  We
also exhibit a proper antipodal orbit for which irreducibility alone does not
force criticality.

Group-invariant kernel energies on compact homogeneous manifolds have also
been used for quadrature and discrepancy \cite{DamelinEtAl2009}.  Those
results optimize an energy over measures or point sets on the homogeneous
space.  Our variable is instead an ambient external direction on a fixed
Frobenius sphere, and the conclusion is a transverse Hessian theorem at the
embedded source orbit.

For half-Grassmannians this proves that the balanced field never becomes
spinodally unstable over $\R$, $\C$, or $\Hh$.  In the complex case a second,
independent Jacobi calculation classifies weak-field block splitting at every
two-level spectrum.  It reveals locally stable unbalanced branches near the
balanced multiplicity.  Hence local Hessians cannot eliminate all
competitors: a complete global proof must compare separated values or prove
a stronger Laplace order.

Our contributions are:
\begin{enumerate}[label=(\roman*)]
\item a positive-mode spherical-Laplacian identity, a necessary criticality
guard, and an all-field Morse--Bott criterion for homogeneous orbital
potentials;
\item strict balanced local rigidity for the real, complex, and quaternionic
half-Grassmann matrix--Bingham models;
\item exact Jacobi--Stein identities and exact constrained-Hessian operators
for every complex two-block multiplicity stratum;
\item the sharp weak-field metastability threshold $k=2r/3$, including its
degenerate sixth-moment boundary; and
\item a coefficient-tail theorem that reduces one global sign problem to
finitely many exact moments at each fixed $(r,k)$.
\end{enumerate}
The addition theorem, matrix-beta law, and Jacobi ensemble are established
ingredients.  Novelty is claimed only for the joined all-field stability and
metastability conclusions, not for those classical tools.

\section{Orbital harmonic positivity}

Let $N\ge3$, let $G$ be a compact group acting orthogonally on $\R^N$, let
$\cO=Gv\subset S^{N-1}$ be an orbit, and let $\mu$ be its invariant
probability measure.  Write $\lambda=(N-2)/2$ and let
\[
 K_\ell(x,y)=\frac{C_\ell^\lambda(\ip{x}{y})}
 {C_\ell^\lambda(1)}                                      \tag{2.1}
\]
be the normalized degree-$\ell$ spherical-harmonic reproducing kernel.
The circle $N=2$ has the analogous Fourier/Chebyshev formulation, but it is
not needed in any application below and is excluded from the stated
Gegenbauer normalization.
For an orthonormal real harmonic basis $\{Y_{\ell j}\}_j$, the addition
theorem gives
\[
 K_\ell(x,y)=c_{N,\ell}\sum_jY_{\ell j}(x)Y_{\ell j}(y),
 \qquad c_{N,\ell}>0.                                     \tag{2.2}
\]

\begin{lemma}[Squared orbital coefficient]\label{lem:square}
For every $\ell\ge0$ and every $v\in\cO$,
\[
 q_\ell:=\int_\cO K_\ell(u,v)\,d\mu(u)
 =c_{N,\ell}\sum_j\left|\int_\cO Y_{\ell j}(u)\,d\mu(u)\right|^2
 \ge0.                                                     \tag{2.3}
\]
If $\cO$ is antipodal, then $q_\ell=0$ for odd $\ell$.
\end{lemma}

\begin{proof}
Invariance makes the first integral independent of the selected $v\in\cO$.
Averaging it once more over $v$ and using (2.2) gives the squared-norm
formula.  Antipodality annihilates every odd harmonic average.
\end{proof}

The evaluation point belonging to the source orbit is essential.  Away from
the orbit the corresponding coefficient is an inner product of two harmonic
vectors and need not be nonnegative.

Let an analytic zonal kernel have the expansion
\[
 \Phi(x)=\sum_{\ell\ge0}a_\ell
 \frac{C_\ell^\lambda(x)}{C_\ell^\lambda(1)},\qquad a_\ell\ge0,              \tag{2.4}
\]
with locally uniform convergence after two derivatives, and define
\[
 U_\Phi(a)=\int_\cO\Phi(\ip{a}{u})\,d\mu(u),\qquad a\in S^{N-1}.             \tag{2.5}
\]

\begin{proposition}[Negative orbital Laplacian]\label{prop:lap}
At every $v\in\cO$,
\[
 \Delta_{S^{N-1}}U_\Phi(v)
 =-\sum_{\ell\ge1}\ell(\ell+N-2)a_\ell q_\ell\le0.        \tag{2.6}
\]
The inequality is strict if $a_\ell q_\ell>0$ for at least one
$\ell\ge1$.
\end{proposition}

\begin{proof}
Integrate (2.4), apply \cref{lem:square}, and use
$\Delta K_\ell=-\ell(\ell+N-2)K_\ell$ in the first argument.
\end{proof}

For $\Phi_t(x)=e^{tx}$ all Gegenbauer coefficients are strictly positive for
$t>0$.  If $\cO$ is a proper orbit, its invariant measure is not uniform on
the ambient sphere.  Since spherical harmonics determine finite measures on
the compact sphere, some nonconstant $q_\ell$ is positive.  Thus
\cref{prop:lap} is strict for every $t>0$.

\section{A critical all-field Morse--Bott theorem}

The potential is $G$-invariant, hence constant along $\cO$, but constancy on
the orbit does not by itself make an orbit point critical: a normal gradient
can remain.  The following elementary guard isolates the missing condition.

\begin{lemma}[Criticality guard]\label{lem:critical}
For every $v\in\cO$, the spherical gradient satisfies
\[
 \nabla_S U_\Phi(v)\in (N_v\cO)^{G_v}.                       \tag{3.1}
\]
Consequently $(N_v\cO)^{G_v}=\{0\}$ implies that the whole orbit is critical.
At a critical orbit, the spherical Hessian annihilates $T_v\cO$, including
all mixed tangent--normal entries.
\end{lemma}

\begin{proof}
Invariance makes the gradient orthogonal to the orbit and fixed by the
stabilizer, proving (3.1).  For an infinitesimal orbit field $X^\#$ one has
$dU_\Phi(X^\#)=0$.  Covariantly differentiating this identity in an arbitrary
spherical direction $Z$ gives
\[
 \operatorname{Hess}_S U_\Phi(Z,X^\#)
 +dU_\Phi(\nabla_ZX^\#)=0.
\]
The second term vanishes at a critical point.  Orbit homogeneity transports
criticality and the Hessian-kernel conclusion from $v$ to every point of
$\cO$.
\end{proof}

\begin{theorem}[Critical homogeneous orbital Morse--Bott stability]\label{thm:mb}
Assume that $\cO\subset S^{N-1}$ is proper and antipodal, that $v$ is a
critical point of $U_\Phi$---for example,
$(N_v\cO)^{G_v}=\{0\}$---and that either
\begin{enumerate}[label=(\alph*)]
\item the real $G_v$-representation on $N_v\cO$ is irreducible; or
\item it is a direct sum of pairwise inequivalent irreducibles and a symmetry
of the sphere fixing $v$ and preserving $U_{\Phi}$ permutes the summands
transitively.
\end{enumerate}
If $\Phi$ satisfies (2.4) with $a_\ell>0$ for every $\ell$, then
$\cO$ is a nondegenerate Morse--Bott manifold of strict local maxima of
$U_\Phi$.  In particular this holds for $\Phi_t(x)=e^{tx}$ at every $t>0$.
\end{theorem}

\begin{proof}
By \cref{lem:critical}, the tangent and mixed Hessian entries vanish.  Schur's
lemma makes the self-adjoint Hessian scalar on every indicated irreducible
normal summand; the extra symmetry makes those scalars equal in case (b).
Thus the spherical Laplacian is the dimension of the normal space times their
common scalar.  It is strictly negative by
\cref{prop:lap}; hence the normal Hessian is negative definite.  Its kernel
is exactly $T_v\cO$, and the Morse--Bott lemma gives transverse strict local
maximality.
\end{proof}

\begin{remark}[Why criticality is essential]\label{rem:torus}
Let $G=SO(2)\times SO(2)$ act on $\R^2\oplus\R^2$ and take
$v=(ae_1,be_1)$ with $a^2+b^2=1$ and $a\ne b$.  The proper antipodal orbit
$aS^1\times bS^1\subset S^3$ has a one-dimensional, hence irreducible,
spherical normal representation, but that representation is trivial.  For
$A(\theta)=(\cos\theta\,e_1,\sin\theta\,e_1)$,
\[
 U_t(A(\theta))=I_0(ta\cos\theta)I_0(tb\sin\theta),
\]
and at $\cos\theta_0=a$, $\sin\theta_0=b$,
\[
 \left.\frac d{d\theta}\log U_t(A(\theta))\right|_{\theta_0}
 =tab\!\left[\frac{I_1(tb^2)}{I_0(tb^2)}-
               \frac{I_1(ta^2)}{I_0(ta^2)}\right].          \tag{3.2}
\]
Its small-field expansion is
$\tfrac12t^2ab(b^2-a^2)+O(t^4)$, so the source point is not critical for all
sufficiently small $t>0$.  This example satisfies the original irreducibility
condition but violates the fixed-vector guard in \cref{lem:critical}.
\end{remark}

\begin{corollary}[Uniform tubular exclusion]\label{cor:tube}
Assume $(N_v\cO)^{G_v}=\{0\}$ and either normal-isotropy condition of
\cref{thm:mb}.  For every compact $[t_0,t_1]\subset(0,\infty)$ there is a tubular
neighborhood $\mathcal T$ of $\cO$ such that
\[
 U_t(a)<U_t(v),\qquad
 a\in\mathcal T\setminus\cO,\quad v\in\cO,\quad t\in[t_0,t_1].           \tag{3.3}
\]
\end{corollary}

\begin{proof}
The smallest magnitude normal-Hessian eigenvalue is continuous on the
compact set $\cO\times[t_0,t_1]$ and is bounded away from zero.  A uniform
Taylor remainder and compactness give the claimed tube.
\end{proof}

\section{Half-Grassmann no-spinodal rigidity}

Let $\mathbb F\in\{\R,\C,\Hh\}$, let $P$ be a Haar rank-$r$ orthogonal
projector on $\mathbb F^{2r}$, and set
\[
 S_P=P-\frac12I,\qquad \rho^2=\norm{S_P}_F^2=\frac r2,\qquad
 u_P=S_P/\rho.                                               \tag{4.1}
\]
The orbit $\cO_{\mathbb F}=\{u_P\}$ is antipodal because $I-P$ has the same
law.  Its real ambient dimensions are
\[
 N_{\R}=r(2r+1)-1,\qquad N_{\C}=4r^2-1,\qquad
 N_{\Hh}=8r^2-2r-1.                                         \tag{4.2}
\]
At
\[
 v_\star=\frac{\operatorname{diag}(I_r,-I_r)}{\sqrt{2r}},                  \tag{4.3}
\]
the off-diagonal block is tangent to the conjugacy orbit.  The spherical
normal space is the direct sum of the two block-diagonal traceless
self-adjoint modules.  Each is irreducible under its stabilizer factor; block
exchange together with antipodality makes their Hessian scalars equal.  More
explicitly, if $W$ swaps the two $r$-blocks, then
$A\mapsto-WAW^*$ fixes $v_\star$, exchanges the normal modules, and preserves
the potential because $\mu$ is both conjugation-invariant and antipodal.
Moreover, a vector fixed by the block stabilizer is scalar on each block;
membership in either traceless normal module forces both scalars to vanish.
Hence $(N_{v_\star}\cO_{\mathbb F})^{G_{v_\star}}=\{0\}$ over
$\mathbb F=\R,\C,\Hh$, and \cref{lem:critical} proves criticality rather than
assuming it.

\begin{theorem}[Balanced no-spinodal theorem]\label{thm:grass}
For every $r\ge2$, every $\mathbb F\in\{\R,\C,\Hh\}$, and every $t>0$, the
balanced field orbit (4.3) is a strict constrained local maximizer of
\[
 A\longmapsto \E\exp\!\bigl(t\rho\ip{A}{u_P}\bigr),
 \qquad \tr A=0,\quad\norm A_F=1.                           \tag{4.4}
\]
Its constrained-Hessian kernel consists exactly of conjugacy-orbit
directions.  Thus the balanced branch has no finite-field radial spinodal.
\end{theorem}

\begin{proof}
The fixed-vector calculation above supplies the criticality hypothesis.
Apply \cref{thm:mb} to the normal decomposition and the block-exchange
symmetry.
\end{proof}

For the complex case let
\[
 z_r(t)=\E e^{t\rho\ip{v_\star}{u_P}}.
\]
The Gegenbauer--Bessel expansion gives the useful scalar certificate
\[
 \boxed{G_r(t):=z_r''(t)+(4r^2-2)\frac{z_r'(t)}t-\frac r2z_r(t)>0.}          \tag{4.5}
\]
Indeed the radial Bessel equation yields
\[
 G_r(t)=\sum_{\ell\ge1}\frac{\ell(\ell+4r^2-3)}{t^2}
 a_{\ell}(t)q_\ell>0.                                      \tag{4.6}
\]
Equivalently, the spherical Laplacian at $v_\star$ equals $-t^2G_r(t)$;
the common normal-Hessian eigenvalue is
\[
 -\frac{t^2G_r(t)}{2(r^2-1)}.                               \tag{4.7}
\]

\section{Complex two-block strata}

We now specialize to $P\sim\Gr_{\C}(r,2r)$.  Put $n=2r$, fix
$1\le k<r$, set $m=n-k$, and let $E_k$ project onto a fixed $k$-plane.
Define
\[
 T=\tr(E_kP),\qquad Y=T-\frac k2,\qquad z_{r,k}(u)=\E e^{uY}.               \tag{5.1}
\]
The $k\times k$ compression $B=E_kPE_k$ has the complex matrix-beta density
\[
 c_{r,k}\det(B)^{r-k}\det(I-B)^{r-k}\mathbf1_{0<B<I}\,dB,                \tag{5.2}
\]
and its eigenvalues form a Jacobi unitary ensemble
\cite{James1964,Muirhead1982,DumitriuEdelman2002}.  The Frobenius-unit
two-block field is
\[
 A_k=c_k\left(E_k-\frac{k}{n}I\right),\qquad
 c_k=\sqrt{\frac{n}{km}},\qquad \tr(A_kP)=c_kY.             \tag{5.3}
\]
If the physical field strength is $s$, write $u=sc_k$.

\begin{lemma}[Two-block criticality]\label{lem:twoblock-critical}
For every $s>0$, the conjugacy orbit of $A_k$ is a critical manifold of
\[
 F_s(A)=\log\E e^{s\tr(AP)}
 \quad\text{on}\quad \{A=A^*: \tr A=0,\ \norm A_F=1\}.
\]
\end{lemma}

\begin{proof}
Under the posterior tilted by $A_k$, the mean $M_s=\E_{A_k}P$ commutes with
the stabilizer $U(k)\times U(m)$.  Hence
$M_s=aE_k+b(I-E_k)$ for scalars $a,b$.  Since $\tr M_s=r$, its traceless part
is proportional to the unique traceless block-scalar direction
$E_k-(k/n)I$, and therefore to $A_k$.  The trace-zero ambient gradient of
$F_s$ is $s(M_s-\tfrac12I)$, so it is radial on the Frobenius sphere at
$A_k$.  Every constrained first variation consequently vanishes.  Unitary
invariance transports criticality along the full conjugacy orbit.
\end{proof}

For the eigenvalues $x_1,\ldots,x_k$ of $B$, set
\[
 S=\sum_{i=1}^k x_i(1-x_i).                                 \tag{5.4}
\]
The reversible complex-Jacobi generator is
\[
 \cL f=\sum_i x_i(1-x_i)\partial_{ii}f+
 \sum_i b_i(x)\partial_if,                                  \tag{5.5}
\]
where
\[
 b_i=(r-k+1)(1-2x_i)+2x_i(1-x_i)
 \sum_{j\ne i}\frac1{x_i-x_j}.                              \tag{5.6}
\]

\begin{lemma}[Jacobi--Stein closure]\label{lem:stein}
With $\Gamma$ the carré du champ of (5.5),
\[
 \cL Y=-2rY,\qquad \Gamma(Y,Y)=S,\qquad
 \cL S=\frac{km}{2}+2Y^2-4rS.                               \tag{5.7}
\]
Under the tilted law $d\mu_u=e^{uY}d\mu/z_{r,k}(u)$,
\[
 2r\,\delta(u)=u\E_uS,\qquad
 2r\E_uY^2=\E_uS+u\E_u(YS),\qquad \delta=(\log z_{r,k})'. \tag{5.8}
\]
\end{lemma}

\begin{proof}
The first two identities follow directly from (5.5).  For the third,
pair the singular interaction terms for $(i,j)$ before simplifying.  The
divided differences of $x(1-x)(1-2x)=x-3x^2+2x^3$ reduce to $T$,
$\sum x_i^2$, and $T^2$; centering cancels every linear term and gives
(5.7).  Reversibility gives
$\E[\cL f\,e^{uY}]=-u\E[\Gamma(f,Y)e^{uY}]$.  Taking $f=Y$ and then
$f=Y^2/2$ gives (5.8).
\end{proof}

\section{Exact constrained Hessians}

Let $q(u)=\E_uY^2=z''/z$.  For $k\ge2$, unitary invariance shows that a
Frobenius-unit traceless split within the positive $k$-block has posterior
variance
\[
 \alpha_+=\frac{k/4-2r\delta/u-q/k}{k^2-1},                  \tag{6.1}
\]
while a unit split within the negative $m$-block has variance
\[
 \alpha_-=\frac{m/4-2r\delta/u-q/m}{m^2-1}.                  \tag{6.2}
\]
For example, (6.1) follows from
\[
 \E_u\left[\tr B^2-\frac{T^2}{k}\right]
 =\frac k4-\E_uS-\frac{q}{k},                               \tag{6.3}
\]
and (6.2) follows by applying the same invariant covariance formula to
the complementary block spectrum
$(1-x_1,\ldots,1-x_k,1^{r-k},0^{r-k})$.

Along a unit tangent geodesic on the Frobenius sphere, the Hessian of the log
partition divided by $s^2$ is posterior variance minus the spherical
Lagrange multiplier.  This is a constrained Hessian because
\cref{lem:twoblock-critical} supplies the required first-order stationarity:
\[
 \ell_{r,k}(u)=\frac{\E_u\tr(A_kP)}s=\frac{n\delta(u)}{kmu}.                 \tag{6.4}
\]
Define
\[
\begin{aligned}
 H_{r,k}^{+}(u)&=z''+\frac{n(nk-1)}m\frac{z'}u-\frac{k^2}{4}z,\\
 H_{r,k}^{-}(u)&=z''+\frac{n(nm-1)}k\frac{z'}u-\frac{m^2}{4}z.
\end{aligned}                                                \tag{6.5}
\]

\begin{theorem}[Two-block Hessian operators]\label{thm:hess}
For every $u>0$, the negative-block identity below holds; when $k\ge2$ the
positive-block identity holds as well:
\[
 \boxed{
 \alpha_+-\ell_{r,k}=-\frac{H_{r,k}^{+}}{k(k^2-1)z},\qquad
 \alpha_--\ell_{r,k}=-\frac{H_{r,k}^{-}}{m(m^2-1)z}.}       \tag{6.6}
\]
For $k=1$ the positive block has no traceless splitting module, so its first
displayed quotient is omitted.  Thus all existing fundamental spectral
splitting signs are reduced exactly to scalar differential inequalities for
one Jacobi trace transform.
\end{theorem}

\begin{proof}
Substitute (5.8), (6.1), and (6.2) into
(6.4); collect terms using $q=z''/z$ and $\delta=z'/z$.
\end{proof}

\section{A sharp metastability threshold}

The centered law of $Y$ is symmetric.  Its second and fourth moments are
\[
 \mu_2=\frac{km}{4(n^2-1)},\qquad
 \mu_4=\frac{3km(km-2)}
 {16(n-3)(n-1)(n+1)(n+3)}.                                  \tag{7.1}
\]
These follow either from the matrix-beta Schur expansion or from the standard
complex projector moments.  Substitution into (6.5) gives the actual
Taylor coefficients
\[
 [u^2]H_{r,k}^{+}
 =\frac{k(4r-3k)(k^2-1)}
 {16(2r-3)(2r-1)(2r+1)(2r+3)}\ge0,                          \tag{7.2}
\]
with equality only at $k=1$, where that splitting module is absent,
and
\[
 [u^2]H_{r,k}^{-}
 =-\frac{(2r-3k)m(m^2-1)}
 {16(2r-3)(2r-1)(2r+1)(2r+3)}.                              \tag{7.3}
\]
When $3k=2r$, (7.3) vanishes.  The exact sixth moment then gives
\[
 [u^4]H_{r,2r/3}^{-}
 =-\frac{4r^2(4r-3)(4r+3)}
 {5832(2r-5)(2r-1)^2(2r+1)^2(2r+5)}<0.                     \tag{7.4}
\]

\begin{theorem}[Weak-field metastability classification]\label{thm:meta}
For sufficiently small nonzero field:
\begin{enumerate}[label=(\roman*)]
\item if $k<2r/3$, the larger block has an unstable traceless split;
\item if $k=2r/3$, the same instability begins at order $u^4$; and
\item if $2r/3<k<r$, both block-splitting Hessian eigenvalues are negative,
so the unbalanced two-block orbit is a genuine local maximum modulo its
conjugacy orbit.
\end{enumerate}
\end{theorem}

\begin{proof}
Use \cref{thm:hess} and the signs in (7.2)--(7.4).  The remaining
off-diagonal directions are conjugacy-orbit zero modes; the two block
modules exhaust the normal spectral directions.
\end{proof}

\begin{table}[H]
\centering
\small
\begin{tabular}{@{}lll@{}}
\toprule
Multiplicity range & Weak-field normal type & All-/high-field information \\
\midrule
$1\le k<2r/3$ & saddle & larger-block split already unstable \\
$k=2r/3$ & saddle at order $u^4$ & quadratic coefficient vanishes \\
$2r/3<k<r$ & metastable local maximum & larger block unstable at high field \\
$k=r$ & strict local maximum & stable for every $u>0$ by \cref{thm:grass} \\
\bottomrule
\end{tabular}
\caption{Exact local classification of the fundamental complex two-block
strata.  ``Metastable'' is a local geometric statement, not a claim of
global phase activity.}
\label{tab:stability}
\end{table}

This theorem rules out the shortcut ``every unbalanced two-block stationary
orbit is a saddle.''  Near-balanced multiplicities are metastable at weak
field.  At high field, however,
\[
 \frac{z''}{z}\longrightarrow\frac{k^2}{4},\qquad
 \frac{z'}{uz}\longrightarrow0,
\]
and hence $H_{r,k}^{-}/z\to(k^2-m^2)/4<0$.  Each such metastable branch
eventually becomes unstable.  This says nothing by itself about which branch
has the largest value.

\section{A rigorous coefficient-tail reduction}

Let $\mu_j=\E Y^j$.  Coefficientwise negativity of $H_{r,k}^{-}$ is
equivalent to
\[
 \frac{\mu_{2j+2}}{\mu_{2j}}
 <\frac{m^2}{4}\frac{2j+1}{2j+1+n(nm-1)/k}.                 \tag{8.1}
\]

\begin{proposition}[Stein moment-ratio bound]\label{prop:tail}
For every $j\ge0$,
\[
 \frac{\mu_{2j+2}}{\mu_{2j}}
 \le\frac{k^2}{4}\frac{2j+1}{2j+1+2rk}.                    \tag{8.2}
\]
Consequently (8.1) holds whenever
\[
 2j(m-k)>m(k^2-1).                                          \tag{8.3}
\]
\end{proposition}

\begin{proof}
The pointwise inequality
\[
 S=\sum_i x_i(1-x_i)\le\frac k4-\frac{Y^2}{k}               \tag{8.4}
\]
follows from Cauchy--Schwarz after writing $x_i=1/2+y_i$.
Reversibility and $\cL Y=-2rY$ give
\[
 2r\E Y^{2j+2}=(2j+1)\E(Y^{2j}S).
\]
Insert (8.4), rearrange, and obtain (8.2).  Direct algebra shows
that its right-hand side is strictly smaller than the right-hand side of
(8.1) precisely under (8.3).
\end{proof}

For every fixed $(r,k)$ this proves the entire sufficiently high-degree tail
of the larger-block Hessian operator.  The remaining number of low degrees
is finite, but it is not uniform as $k\uparrow r$.  Finite numerical screens
therefore do not constitute an all-rank proof.

\section{Phase implications and scope}

The balanced branch is a strict local maximum at every field by
\cref{thm:grass}.  Hence any hypothetical competitor that overtakes it must
be separated from the balanced orbit.  It cannot emerge continuously through
a vanishing balanced Hessian.  On a compact field interval,
\cref{cor:tube} removes a uniform tube around the balanced orbit from any
global search.

The metastability theorem adds a complementary warning.  Several unbalanced
two-block strata are also local maxima at weak field, so no argument based
only on stationary-point Hessians can establish the global envelope.  A
global theorem must instead supply at least one of:
\begin{enumerate}[label=(\alph*)]
\item an all-field Laplace order between external spectra;
\item a value comparison for every two-block multiplicity plus a theorem
reducing global maximizers to two blocks; or
\item a nonlocal positivity or total-positivity certificate for the Jacobi
Hankel determinants.
\end{enumerate}
The trace transform in (5.1) is connected to differential recurrences
for the Jacobi ensemble \cite{ForresterKumar2022}; those recurrences do not,
by themselves, provide the needed fixed-Frobenius ordering.

The paper proves local orbital geometry and exact metastability.  It does
\emph{not} prove global maximization by the balanced field, a complete
finite-field phase diagram, or an all-distortion Shannon rate--distortion
function.  It also does not interpret local maxima as physical thermodynamic
states unless a separate model supplies such an interpretation.

\section{Reproducibility}

The accompanying lightweight replay uses exact rational arithmetic and
symbolic factorization to verify (7.1)--(7.4), the algebraic
equivalence in (8.3), and the ambient dimensions in (4.2).  It
also evaluates the noncritical torus counterexample in \cref{rem:torus} and
checks the positive Bessel-mode formula on a bounded grid.  The replay
is evidence against transcription errors, not a computer-assisted proof of
the harmonic or Jacobi theorems; all logical steps used above are given in
the text.

\section{Conclusion}

Positive-definite orbital potentials have a rigid Laplacian identity at their
source orbit: their nonconstant harmonic content enters with one sign.  At a
critical source orbit, an isotropy condition upgrades that trace identity to
a complete all-field Morse--Bott maximum theorem; without criticality this is
false, as the explicit torus orbit shows.  For half-Grassmann
matrix--Bingham models the normal fixed-vector space vanishes, so the repaired
criterion applies and
excludes a balanced spinodal over all three classical division algebras.
The exact Jacobi reduction then exposes a less obvious structure in the
complex problem: near-balanced unbalanced spectra are weak-field metastable,
while more asymmetric spectra are saddles.  These results sharply constrain
any unresolved intermediate phase while preserving the central global gate
rather than hiding it behind finite moments or numerical evidence.

\appendix
\section{Taylor normalization at the degenerate threshold}

For completeness, with $z(u)=\sum_{j\ge0}\mu_{2j}u^{2j}/(2j)!$, an operator
$H=z''+Az'/u-cz$ has
\[
 [u^2]H=\frac12\left[\mu_4\left(1+\frac A3\right)-c\mu_2\right],\qquad
 [u^4]H=\frac1{24}\left[\mu_6\left(1+\frac A5\right)-c\mu_4\right].       \tag{A.1}
\]
At $k=2r/3$, $m=4r/3$,
\[
 \mu_4=\frac{r^2}{27(2r-1)(2r+1)},                          \tag{A.2}
\]
and the Schur expansion of the complex matrix-beta law gives
\[
 \mu_6=\frac{10r^2(4r^4-27r^2+9)}
 {243(2r-5)(2r-1)^2(2r+1)^2(2r+5)}.                         \tag{A.3}
\]
Substitution of $A=n(nm-1)/k$ and $c=m^2/4$ in (A.1) yields
(7.4).  The factorials in (A.1) are included explicitly because
omitting them leaves the signs unchanged but doubles the quadratic
coefficient and multiplies the quartic coefficient by $24$.

\section{Algebra behind the Jacobi closure}

We record the drift calculation in \cref{lem:stein}.  Write
$p_1=T=\sum_i x_i$ and $p_2=\sum_i x_i^2$, so $S=p_1-p_2$.  Pairwise
symmetrization of the interaction term gives
\[
 4\sum_{i<j}\frac{x_i^2(1-x_i)-x_j^2(1-x_j)}{x_i-x_j}
 =4(k-1)T-(4k-6)p_2-2T^2.                                  \tag{B.1}
\]
Adding the one-particle drift and diffusion terms yields
\[
 \cL p_1=kr-2rT,
 \qquad
 \cL p_2=(2r+2k)T-4rp_2-2T^2.                              \tag{B.2}
\]
Therefore
\[
 \cL S=kr-2kT+2T^2-4rS
 =\frac{k(2r-k)}2+2\left(T-\frac k2\right)^2-4rS,           \tag{B.3}
\]
which is (5.7).  The matrix-beta moment identity used in the replay is
\[
 \E s_\lambda(B)=s_\lambda(1^k)
 \frac{[r]_\lambda}{[2r]_\lambda},                         \tag{B.4}
\]
with the complex generalized rising factorial
$[a]_\lambda=\prod_i(a-i+1)_{\lambda_i}$.  Expanding
$(\tr B)^j=\sum_{\lambda\vdash j}f^\lambda s_\lambda(B)$ and centering gives
(7.1), (A.2), and (A.3) by exact rational simplification.

\printbibliography
\end{document}
